\clearpage
\VerdierPhysicalPageAnchor{0073}%
\VerdierSemanticAnchor{ch1-2-3-4}
\noindent\textbf{2.3.4.}\quad
取 $\varphi$ 为唯一的映射 $[n]\xrightarrow{\varphi}[1]$，并令群
$\mathbb Z[n]$ 通过
$\varphi\colon\mathbb Z[n]\longrightarrow\mathbb Z$ 作用在
$\operatorname{comp}(C)$ 上。根据 (1.5.2)，此时存在唯一的由
Koszul 余循环扭曲的 $\mathbb Z[n]$-函子：

\VerdierSemanticAnchor{ch1-2-3-4-eq1}
\begin{equation}
  \left(\int_n^\Sigma,p_n\right)\colon
  \operatorname{comp}^{n}(C)\longrightarrow\operatorname{comp}(C)
  \tag{2.3.4.1}
\end{equation}

\VerdierSemanticAnchor{ch1-2-3-4-characterization}
使得：
\begin{equation*}
  p_n(e_i)=p_{i,\varphi}.
\end{equation*}

\VerdierSemanticAnchor{ch1-2-3-4-notation}
今后，记号 $\int_n^\Sigma$ 表示函子
$\int_n^\Sigma$ 连同上述扭曲 $\mathbb Z[n]$-函子结构。

\VerdierSemanticAnchor{ch1-2-3-5}
\noindent\textbf{2.3.5.}\quad
在 (2.3.2)、(2.3.3) 和 (2.3.4) 中展开的讨论，经过显然的修改后，
同样适用于函子 $\int_\varphi^\Pi$。

\VerdierSemanticAnchor{ch1-2-3-5-summability-restriction}
当可数直和与直积在 $C$ 中不一定可表示时，只要将范围限制在具有适当
可求和性的复形上，上述结果仍适用于函子 $\int_\varphi$。

\VerdierSemanticAnchor{ch1-2-4}
\subsection{函子到复形的延拓}

\VerdierSemanticAnchor{ch1-2-4-1}
\noindent\textbf{2.4.1.}\quad
存在一个典范函子：

\VerdierSemanticAnchor{ch1-2-4-1-eq1}
\begin{equation}
  \iota_n\colon
  \operatorname{comp}^{n}(C)^\circ
  \longrightarrow
  \operatorname{comp}^{n}(C^\circ)
  \tag{2.4.1.1}
\end{equation}
它由下列公式定义：

\VerdierSemanticAnchor{ch1-2-4-1-object-formula}
\begin{equation*}
  (\iota_nX)^\sigma=X^{-\sigma},
  \qquad X\in\operatorname{Ob}\bigl(\operatorname{comp}^{n}(C)\bigr),
\end{equation*}

\VerdierSemanticAnchor{ch1-2-4-1-eq2}
\begin{equation}
  d_{\iota_nX}^{\sigma,i}=d_X^{-\sigma-e_i,i},
  \tag{2.4.1.2}
\end{equation}

\VerdierSemanticAnchor{ch1-2-4-1-morphism-formula}
\begin{equation*}
  (\iota_nf)^\sigma=f^{-\sigma},
  \qquad f\in\operatorname{Fl}\bigl(\operatorname{comp}^{n}(C)\bigr).
\end{equation*}

\VerdierSemanticAnchor{ch1-2-4-1-graded-isomorphism}
函子 $\iota_n$ 是 $\mathbb Z[n]$-范畴的一个同构
（$\operatorname{comp}^{n}(C^\circ)$ 的 $\mathbb Z[n]$-范畴结构由 (2.3.1) 定义；
$\operatorname{comp}^{n}(C)^\circ$ 的分次结构由 (2.3.1) 和 (1.3.9) 定义）。

\VerdierSemanticAnchor{ch1-2-4-1-sigma-opposite-setup}
假定可数直和在 $C$ 中可表示，并设
$\varphi\colon[n]\longrightarrow[m]$ 为一个映射。记：

\VerdierSemanticAnchor{ch1-2-4-1-eq3}
\begin{equation}
  \left(\int_{\varphi,C}^{\Sigma}\right)^\circ\colon
  \operatorname{comp}^{n}(C)^\circ
  \longrightarrow
  \operatorname{comp}^{m}(C)^\circ,
  \tag{2.4.1.3}
\end{equation}
